% Shared macros for the Lake Country fishing guide.
%
% Loaded after common/preamble.tex by every leaf in this project. Everything
% here exists to keep three families of page consistent: species cards, rig
% cards, and lake maps.

% ---------------------------------------------------------------------------
% Species card header. The plate sits full width under the masthead, because a
% fish is identified by its whole profile and a thumbnail defeats the purpose.
%   \speciesplate{file}{caption}
%
% The Makefile builds each leaf from its own directory with src/ on TEXINPUTS,
% so shared art is named from the project root rather than relatively.
% ---------------------------------------------------------------------------
\newcommand{\telosplateroot}{lake-country-fishing/claude/shared/plates}
\newcommand{\speciesplate}[2]{%
  \begingroup
  \centering
  \includegraphics[width=0.86\linewidth]{\telosplateroot/#1.png}\par
  \vspace{-0.1em}
  \begin{minipage}{0.86\linewidth}
    \rule{\linewidth}{0.35pt}\par
    \vspace{0.08em}
    {\fontsize{7}{8}\selectfont #2\par}
  \end{minipage}\par
  \endgroup
  \vspace{0.15em}
}

% ---------------------------------------------------------------------------
% Species page header: plate on the left, identification facts on the right.
%
% A species sheet is a one-pager, and that is a hard contract -- it has to be
% usable as a single sheet in a tackle box. Stacking a full-width plate above
% the fact strip costs about four centimetres of height that the page does not
% have, so the two sit side by side instead. The body of the environment is a
% run of \telosfact rows.
%
%   \begin{speciestop}{plate-slug}{plate credit}
%     \telosfact{...}{...}
%   \end{speciestop}
% ---------------------------------------------------------------------------
\newenvironment{speciestop}[2]
  {\par\noindent
   \begin{minipage}[t]{0.395\linewidth}
     \vspace{0pt}%
     \includegraphics[width=\linewidth]{\telosplateroot/#1.png}\par
     \vspace{0.15em}
     \rule{\linewidth}{0.35pt}\par
     \vspace{0.1em}
     {\fontsize{6.2}{7.2}\selectfont #2\par}
   \end{minipage}%
   \hfill
   \begin{minipage}[t]{0.575\linewidth}
     \vspace{0pt}\small
     \setlength{\parskip}{0pt}%
     \renewcommand{\arraystretch}{1.22}%
     \begin{tabular}{@{}>{\bfseries\raggedright\arraybackslash}p{0.21\linewidth}@{\hspace{0.55em}}>{\raggedright\arraybackslash}p{\dimexpr0.79\linewidth-0.55em\relax}@{}}}
  {\end{tabular}%
   \end{minipage}\par
   \vspace{0.45em}}

% Standard credit line for the Denton plates. Every species page uses it, so
% the wording of the public-domain claim cannot drift between pages.
\newcommand{\dentoncredit}[1]{%
  #1 --- chromolithograph by S.\,F.\ Denton, New York Commissioners of
  Fisheries, Game and Forests, 1896--1902. Public domain; converted to
  grayscale for printing.}
\newcommand{\fmibcredit}[1]{%
  #1 --- plate from the Freshwater and Marine Image Bank, University of
  Washington. Public domain; converted to grayscale for printing.}

% ---------------------------------------------------------------------------
% Where-and-when block. Every species page answers the same four questions in
% the same order and the same place on the page, so the guide can be used by
% flipping rather than reading.
% ---------------------------------------------------------------------------
\newenvironment{whenwhere}
  {\begingroup\small
   \setlength{\parskip}{0pt}%
   \renewcommand{\arraystretch}{1.3}%
   \begin{tabular}{@{}>{\bfseries\raggedright\arraybackslash}p{0.13\linewidth}@{\hspace{0.6em}}>{\raggedright\arraybackslash}p{\dimexpr0.87\linewidth-0.6em\relax}@{}}}
  {\end{tabular}\par\endgroup}

% ---------------------------------------------------------------------------
% Map primitives.
%
% Both lake maps are drawn in a normalized frame where one unit is one
% centimetre on the page. Depth is encoded by line weight alone: the shoreline
% is heaviest, and each successive contour is lighter, so the map reads on a
% black-only printer exactly as it does on screen.
% ---------------------------------------------------------------------------
\tikzset{
  shore/.style={draw=telosink,line width=.82pt,decorate,
    decoration={random steps,segment length=2.2pt,amplitude=.14pt},
    preaction={draw=black!17,line width=2.2pt},
    postaction={draw=black!58,line width=.24pt}},
  contour5/.style={draw=telosink,line width=0.54pt,
    preaction={draw=black!14,line width=1.02pt}},
  contour10/.style={draw=telosink,line width=0.46pt,
    preaction={draw=black!13,line width=.9pt}},
  contour20/.style={draw=black!82,line width=0.42pt,
    preaction={draw=black!12,line width=.82pt}},
  contour30/.style={draw=black!82,line width=0.36pt,dash pattern=on 3pt off 1.4pt,
    preaction={draw=black!11,line width=.72pt}},
  contour40/.style={draw=black!78,line width=0.38pt,
    preaction={draw=black!10,line width=.72pt}},
  contour80/.style={draw=black!78,line width=0.34pt,dash pattern=on 2pt off 1.2pt,
    preaction={draw=black!10,line width=.66pt}},
  weedbed/.style={draw=telosmute,line width=0.35pt,dash pattern=on 1.2pt off 1pt},
  % The shallow shelf is stippled rather than shaded: a dot screen survives a
  % black-only print driver, where a light grey fill turns into either nothing
  % or a solid block depending on the printer.
  telosshallow/.style={pattern=crosshatch dots,pattern color=black!48},
  maplabel/.style={font=\fontsize{5.6}{6.4}\selectfont,inner sep=0.8pt},
  depthlabel/.style={font=\fontsize{5}{5.6}\selectfont\itshape,text=telosmute,inner sep=0.6pt},
}

% A numbered fishing spot. The disc is solid black with reversed type so it
% stays findable against contour lines at any print density.
\newcommand{\spot}[3]{%
  \node[circle,fill=telosink,text=telospaper,inner sep=0pt,minimum size=3.1mm,
        font=\fontsize{6}{6}\selectfont\bfseries] at (#1,#2) {#3};}

% A boat landing marker.
\newcommand{\landing}[2]{%
  \node[rectangle,draw=telosink,fill=telospaper,line width=0.5pt,inner sep=0pt,
        minimum size=2.6mm,font=\fontsize{5.5}{5.5}\selectfont\bfseries] at (#1,#2) {L};}

% North arrow, placed by its tail.
\newcommand{\northarrow}[2]{%
  \begin{scope}[shift={(#1,#2)}]
    \draw[telos line,-{Stealth[length=1.6mm]}] (0,0) -- (0,0.55);
    \node[font=\fontsize{6}{6}\selectfont\bfseries,anchor=south] at (0,0.55) {N};
  \end{scope}}

% Scale bar, given a length in cm and its label.
\newcommand{\scalebarof}[4]{%
  \begin{scope}[shift={(#1,#2)}]
    \draw[telos hair] (0,0) -- (#3,0);
    \draw[telos hair] (0,-0.06) -- (0,0.06);
    \draw[telos hair] (#3,-0.06) -- (#3,0.06);
    \node[font=\fontsize{5.5}{5.5}\selectfont,anchor=north] at (#3/2,-0.04) {#4};
  \end{scope}}

% Map key entry: a short line of the same weight as the contour it explains.
\newcommand{\mapkeyline}[2]{%
  \raisebox{0.18ex}{\tikz[baseline=0pt]{\draw[#1] (0,0) -- (0.5,0);}}~{\fontsize{6}{7}\selectfont #2}}

% Map key entry for the stippled shallow shelf.
\newcommand{\mapkeystipple}[1]{%
  \raisebox{-0.35ex}{\tikz[baseline=0pt]{
    \fill[telosshallow] (0,0) rectangle (0.5,0.28);
    \draw[telos hair] (0,0) rectangle (0.5,0.28);}}~{\fontsize{6}{7}\selectfont #1}}

% ---------------------------------------------------------------------------
% Rig diagrams. A rig is drawn against a vertical water column: surface rule at
% the top, bottom hatching at the base, everything between drawn to scale in
% the same units so two rigs can be compared by eye.
% ---------------------------------------------------------------------------
\newcommand{\watersurface}[3]{%
  \draw[telos hair] (#1,#3) -- (#2,#3);
  \draw[telos water] (#1,#3-0.08) -- (#2,#3-0.08);}

\newcommand{\lakebottom}[3]{%
  \draw[telos line] (#1,#3) -- (#2,#3);
  \foreach \x in {#1,...,#2} {
    \draw[telos hair] (\x,#3) -- (\x-0.18,#3-0.18);}}

% Terminal-tackle glyphs, all drawn from a single anchor so a rig can be
% assembled by placing anchors down a line.
\newcommand{\rigswivel}[2]{%
  \draw[telos line] (#1,#2) circle (0.07);
  \draw[telos line] (#1,#2+0.07) -- (#1,#2+0.16);
  \draw[telos line] (#1,#2-0.07) -- (#1,#2-0.16);}

\newcommand{\rigsplitshot}[2]{\fill[telosink] (#1,#2) circle (0.075);}

\newcommand{\rigbead}[2]{%
  \draw[telos line] (#1,#2) circle (0.08);
  \draw[telos hair] (#1-0.08,#2) -- (#1+0.08,#2);}

% An egg or walking sinker: a filled lozenge the line runs through.
\newcommand{\rigsinker}[2]{%
  \fill[telosink] (#1,#2) ellipse (0.1 and 0.17);
  \draw[telospaper,line width=0.4pt] (#1,#2-0.2) -- (#1,#2+0.2);}

% A bank or bell sinker hanging from a dropper.
\newcommand{\rigbell}[2]{%
  \fill[telosink] (#1-0.09,#2) -- (#1+0.09,#2) -- (#1+0.05,#2-0.26) -- (#1-0.05,#2-0.26) -- cycle;}

% A hook, drawn hanging from its eye at the anchor.
\newcommand{\righook}[2]{%
  \draw[telos line] (#1,#2) -- (#1,#2-0.3);
  \draw[telos line] (#1,#2-0.3) arc (180:-90:0.11);
  \draw[telos line] (#1+0.22,#2-0.19) -- (#1+0.17,#2-0.09);}

% A jig head: a ball at the bend with the eye at 90 degrees.
\newcommand{\rigjig}[2]{%
  \fill[telosink] (#1,#2-0.28) circle (0.11);
  \draw[telos line] (#1,#2) -- (#1,#2-0.2);
  \draw[telos line] (#1+0.09,#2-0.3) arc (0:-120:0.16);}

% A float. Filled lower half so it reads as sitting in the water.
\newcommand{\rigfloat}[2]{%
  \draw[telos line] (#1,#2) ellipse (0.13 and 0.26);
  \begin{scope}
    \clip (#1-0.15,#2-0.28) rectangle (#1+0.15,#2);
    \fill[telosink] (#1,#2) ellipse (0.13 and 0.26);
  \end{scope}
  \draw[telos line] (#1,#2+0.26) -- (#1,#2+0.4);}

% A bobber stop knot on the line.
\newcommand{\rigstop}[2]{%
  \draw[telos line] (#1-0.07,#2-0.05) -- (#1+0.07,#2+0.05);
  \draw[telos line] (#1-0.07,#2+0.05) -- (#1+0.07,#2-0.05);}

% A soft-plastic worm hanging from a hook, drawn as a tapered body.
\newcommand{\rigworm}[3]{%
  \draw[telos line] (#1,#2) .. controls (#1+0.28,#2-0.25) and (#1-0.24,#2-0.5) ..
    (#1+0.06,#2-#3);}

% A minnow or shiner, side on, hanging from an anchor point.
\newcommand{\rigminnow}[2]{%
  \draw[telos line] (#1,#2) .. controls (#1+0.42,#2+0.16) and (#1+0.42,#2-0.16) .. (#1,#2)
    -- cycle;
  \draw[telos line] (#1+0.42,#2) -- (#1+0.62,#2+0.13) -- (#1+0.62,#2-0.13) -- cycle;
  \fill[telosink] (#1+0.09,#2+0.03) circle (0.025);}

% Callout: a hairline leader from a component to a label set clear of the rig.
%   \riglabel{from x}{from y}{to x}{to y}{anchor}{text}
\newcommand{\riglabel}[6]{%
  \draw[telos hair] (#1,#2) -- (#3,#4);
  \node[maplabel,anchor=#5,font=\fontsize{6.4}{7.4}\selectfont] at (#3,#4) {#6};}

% Depth annotation down the left edge of a rig diagram.
\newcommand{\rigdepth}[4]{%
  \draw[telos dim] (#1,#2) -- (#1,#3);
  \node[maplabel,anchor=east,font=\fontsize{6.4}{7.4}\selectfont,rotate=90]
    at (#1-0.12,{(#2+#3)/2}) {#4};}

% The standard water column a rig is drawn in: surface rule at the top,
% hatched bottom at zero. Every rig sheet uses the same scale so two rigs can be
% compared by holding the sheets side by side.
\newcommand{\rigframe}[3]{%
  \watersurface{0}{#1}{#2}
  \lakebottom{0}{#1}{0}
  \node[maplabel,anchor=west,font=\fontsize{6}{7}\selectfont] at (0.05,#2+0.22) {surface};
  \node[maplabel,anchor=west,font=\fontsize{6}{7}\selectfont] at (0.05,-0.32) {#3};}

% ---------------------------------------------------------------------------
% Cross-reference helper. The guide is a set of loose printed sheets, so a
% reference names the sheet rather than a page number.
% ---------------------------------------------------------------------------
% \sheetref now lives in common/preamble.tex, shared with the other projects.

% ---------------------------------------------------------------------------
% Regulation strip.
%
% Deliberately not a \teloshazard. That box means "this can hurt you" and it is
% reserved for the electricity curriculum; spending it on a bag limit would
% train the reader to ignore it. A rule limit is a legal fact, so it gets a
% compact ruled strip at the foot of the sheet instead -- which also buys back
% the two centimetres that keep a species sheet on one page.
% ---------------------------------------------------------------------------
\newcommand{\regsstrip}[1]{%
  \par\vspace{0.35em}
  \noindent
  \begin{minipage}{\linewidth}
    \rule{\linewidth}{0.9pt}\par
    \vspace{0.3em}
    {\footnotesize\textbf{Regulations on these waters, 2026--27.} #1
     \emph{These change; only the Wisconsin DNR's current publication for this
     exact water body is authoritative.}\par}
    \vspace{0.25em}
    \rule{\linewidth}{0.3pt}
  \end{minipage}\par}

% Field proof belongs beside advice: a reader should know what observation
% would confirm the identification, and what result would disprove it.
\newcommand{\speciesfieldproof}[4]{%
  \section{Prove it in hand}
  \begin{center}
  \begin{tikzpicture}[x=1cm,y=1cm]
    \foreach \x/\n/\t in {0/1/observe,2.75/2/compare,5.5/3/measure,8.25/4/decide}{
      \node[circle,draw=telosink,line width=.7pt,minimum size=5mm,
            font=\scriptsize\bfseries] at (\x,0) {\n};
      \node[font=\scriptsize\bfseries,anchor=west] at (\x+.34,0) {\t};}
    \foreach \x in {1.05,3.8,6.55}{\draw[-{Stealth[length=1.6mm]},telos hair] (\x,0)--(\x+1.08,0);}
  \end{tikzpicture}
  \end{center}
  \begin{telosfacts}[0.18]
    \telosfact{Ask first}{#1}
    \telosfact{Rule out}{#2}
    \telosfact{Measure}{Close the mouth, pinch the tail only if the rule says to, and read total length against a rigid board: #3}
    \telosfact{Field proof}{#4 Photograph the whole side profile and the decisive feature before release if the identification remains uncertain.}
  \end{telosfacts}
}

% A rig is finished only after it passes a dry test and a water test.
\newcommand{\rigfieldproof}[4]{%
  \section{Prove the rig}
  \begin{center}
  \begin{tikzpicture}[x=1cm,y=1cm,node distance=.35cm]
    \node[draw,minimum width=2.75cm,minimum height=.72cm,align=center,font=\scriptsize] (a)
      {\textbf{1 · bench}\\pull every knot};
    \node[draw,minimum width=2.75cm,minimum height=.72cm,align=center,font=\scriptsize,right=of a] (b)
      {\textbf{2 · shallows}\\watch it work};
    \node[draw,minimum width=2.75cm,minimum height=.72cm,align=center,font=\scriptsize,right=of b] (c)
      {\textbf{3 · fishing depth}\\confirm contact};
    \draw[-{Stealth[length=1.8mm]}] (a)--(b);
    \draw[-{Stealth[length=1.8mm]}] (b)--(c);
  \end{tikzpicture}
  \end{center}
  \begin{telosfacts}[0.19]
    \telosfact{Bench test}{#1}
    \telosfact{Water test}{#2}
    \telosfact{Proof at depth}{#3}
    \telosfact{If it fails}{#4 Change one variable, repeat the same cast or drop, and keep the first setting that restores the expected proof.}
  \end{telosfacts}
  \section{Record the adjustment}
  \begin{tabularx}{\linewidth}{@{}c p{.16\linewidth}p{.15\linewidth}p{.17\linewidth}p{.16\linewidth}L@{}}
  \toprule
  \textbf{Try}&\textbf{Depth / distance}&\textbf{Sink count}&
  \textbf{Observed action}&\textbf{Fish response}&\textbf{Next single change}\\
  \midrule
  1&&&&&\\[.55em]
  2&&&&&\\[.55em]
  3&&&&&\\
  \bottomrule
  \end{tabularx}

  \begin{telospanel}{Questions before another cast}
  Did the rig reach the intended depth? What visible line, rod-tip, bottom, or
  lure response proves that? Did the connection and hook point survive the
  retrieve? Was a follow, bite, snag, or refusal observed rather than assumed?
  Name one variable to change and one to hold constant. If the observation
  cannot answer those questions, repeat the same short controlled cast before
  changing the setup.
  \end{telospanel}

  \begin{center}
  \begin{tikzpicture}[x=1cm,y=1cm]
    \draw[telos hair] (0,0) rectangle (10.8,3.2);
    \foreach \x in {.6,1.2,...,10.2}\draw[telos hair] (\x,0)--(\x,3.2);
    \foreach \y in {.6,1.2,...,3.0}\draw[telos hair] (0,\y)--(10.8,\y);
    \node[anchor=south west,font=\scriptsize\bfseries,fill=white]
      at (.15,2.88) {Sketch the cast: surface, cover, target depth, line angle, and where the response occurred};
  \end{tikzpicture}
  \end{center}

  \begin{tabularx}{\linewidth}{@{}p{.25\linewidth}L@{}}
  \textbf{Evidence-based verdict} &
  \fbox{\phantom{x}} repeat unchanged \quad
  \fbox{\phantom{x}} alter one variable \quad
  \fbox{\phantom{x}} retire or repair the rig\\[.8em]
  \textbf{Reason and next prediction} & \hrulefill\\[1.05em]
  & \hrulefill
  \end{tabularx}
}

% Terminal-tackle drawing library, used by every rig sheet.
% Terminal-tackle drawing library.
%
% The species plates in this guide are public-domain chromolithographs and they
% set the standard the rest of the guide has to meet. No equivalent exists for
% terminal tackle: a drop-shot rig and a wacky rig postdate the era of the plates
% by a century, so there is nothing to borrow. These components are therefore
% drawn to match the plates rather than to be schematic -- correct proportions,
% a heavier outline than a circuit diagram would use, and enough interior detail
% that a reader can recognise the real object in their hand.
%
% Every component is drawn from a single anchor with an explicit scale, so a rig
% is assembled by placing anchors down a line. One unit is one centimetre at
% scale 1, which is roughly life size for a size 1/0 hook.

% ---------------------------------------------------------------------------
% Line and leader
% ---------------------------------------------------------------------------
\tikzset{
  fishline/.style={draw=telosink,line width=0.55pt},
  fishleader/.style={draw=telosink,line width=0.45pt},
  fishheavy/.style={draw=telosink,line width=0.9pt},
  fishdetail/.style={draw=telosink,line width=0.3pt},
  fishwire/.style={draw=telosink,line width=0.55pt,dash pattern=on 1.6pt off 1.1pt},
  riglabel/.style={font=\fontsize{6.6}{7.6}\selectfont},
  rigtag/.style={font=\fontsize{6.6}{7.6}\selectfont\itshape},
}

% ---------------------------------------------------------------------------
% Hooks
% ---------------------------------------------------------------------------

% A single J hook hanging from its eye. \hookJ{x}{y}{scale}
\newcommand{\hookJ}[3]{%
  \begin{scope}[shift={(#1,#2)},scale=#3]
    \draw[fishheavy] (0,-0.10) circle (0.10);
    \draw[fishheavy] (0,-0.20) -- (0,-0.92);
    \draw[fishheavy] (0,-0.92) .. controls (0,-1.20) and (0.22,-1.34) ..
      (0.40,-1.24) .. controls (0.54,-1.16) and (0.56,-1.00) .. (0.50,-0.86);
    \draw[fishheavy] (0.50,-0.86) -- (0.44,-0.50);
    \draw[fishdetail] (0.455,-0.62) -- (0.35,-0.70);
  \end{scope}}

% A wide-gap offset worm hook -- the Texas-rig hook, with its characteristic
% double bend just below the eye. \hookOffset{x}{y}{scale}
\newcommand{\hookOffset}[3]{%
  \begin{scope}[shift={(#1,#2)},scale=#3]
    \draw[fishheavy] (0,-0.10) circle (0.10);
    \draw[fishheavy] (0,-0.20) -- (0,-0.34)
      -- (0.16,-0.44) -- (0.02,-0.56) -- (0.02,-1.00);
    \draw[fishheavy] (0.02,-1.00) .. controls (0.02,-1.34) and (0.32,-1.50) ..
      (0.54,-1.36) .. controls (0.70,-1.26) and (0.70,-1.06) .. (0.62,-0.90);
    \draw[fishheavy] (0.62,-0.90) -- (0.54,-0.46);
    \draw[fishdetail] (0.565,-0.60) -- (0.45,-0.68);
  \end{scope}}

% A treble hook. \hookTreble{x}{y}{scale}
\newcommand{\hookTreble}[3]{%
  \begin{scope}[shift={(#1,#2)},scale=#3]
    \draw[fishheavy] (0,-0.10) circle (0.10);
    \draw[fishheavy] (0,-0.20) -- (0,-0.80);
    \foreach \s in {-1,1} {
      \draw[fishheavy] (0,-0.80) .. controls (0,-1.06) and (\s*0.22,-1.18) ..
        (\s*0.38,-1.06) .. controls (\s*0.50,-0.96) and (\s*0.50,-0.82) ..
        (\s*0.44,-0.70);
      \draw[fishheavy] (\s*0.44,-0.70) -- (\s*0.39,-0.38);
      \draw[fishdetail] (\s*0.405,-0.48) -- (\s*0.30,-0.545);}
    \draw[fishheavy] (0,-0.80) .. controls (0.02,-1.14) and (0.12,-1.24) .. (0.16,-1.12)
      .. controls (0.20,-1.00) and (0.14,-0.92) .. (0.06,-0.86);
  \end{scope}}

% A jig head: lead ball moulded at the bend, eye at ninety degrees to the shank.
\newcommand{\jighead}[3]{%
  \begin{scope}[shift={(#1,#2)},scale=#3]
    \draw[fishheavy] (0,0) circle (0.09);
    \draw[fishheavy] (0,-0.09) -- (0,-0.30);
    \fill[telosink] (0,-0.52) circle (0.24);
    \draw[telospaper,line width=0.5pt] (-0.08,-0.44) arc (150:210:0.16);
    \draw[fishheavy] (0.16,-0.64) -- (0.42,-0.94);
    \draw[fishheavy] (0.42,-0.94) .. controls (0.60,-1.12) and (0.60,-0.86) ..
      (0.46,-0.70);
    \draw[fishheavy] (0.46,-0.70) -- (0.40,-0.42);
    \draw[fishdetail] (0.415,-0.52) -- (0.31,-0.58);
  \end{scope}}

% ---------------------------------------------------------------------------
% Weights
% ---------------------------------------------------------------------------

% Bullet (cone) weight, nose up, line running through. \wtBullet{x}{y}{scale}
\newcommand{\wtBullet}[3]{%
  \begin{scope}[shift={(#1,#2)},scale=#3]
    \fill[telosink] (0,0) .. controls (0.20,-0.14) and (0.26,-0.40) .. (0.26,-0.62)
      -- (-0.26,-0.62) .. controls (-0.26,-0.40) and (-0.20,-0.14) .. (0,0) -- cycle;
    \draw[telospaper,line width=0.55pt] (0,0.06) -- (0,-0.68);
    \draw[fishdetail] (-0.26,-0.62) -- (0.26,-0.62);
  \end{scope}}

% Egg sinker. \wtEgg{x}{y}{scale}
\newcommand{\wtEgg}[3]{%
  \begin{scope}[shift={(#1,#2)},scale=#3]
    \fill[telosink] (0,0) ellipse (0.22 and 0.34);
    \draw[telospaper,line width=0.5pt] (0,0.34) -- (0,0.22);
    \draw[telospaper,line width=0.5pt] (0,-0.34) -- (0,-0.22);
    \draw[fishheavy] (0,0.34) -- (0,0.52);
    \draw[fishheavy] (0,-0.34) -- (0,-0.52);
  \end{scope}}

% Walking sinker: the banana-shaped bottom-dragging weight of a lindy rig.
\newcommand{\wtWalking}[3]{%
  \begin{scope}[shift={(#1,#2)},scale=#3]
    \fill[telosink] (-0.06,0.04) .. controls (0.24,-0.06) and (0.40,-0.44) ..
      (0.34,-0.78) .. controls (0.18,-0.72) and (0.02,-0.40) .. (-0.20,-0.10)
      .. controls (-0.16,-0.02) and (-0.12,0.02) .. (-0.06,0.04) -- cycle;
    \draw[fishheavy] (-0.10,0.06) -- (-0.14,0.30);
    \draw[fishheavy] (-0.14,0.30) circle (0.075);
  \end{scope}}

% Bell (bank) sinker hanging from an eye. \wtBell{x}{y}{scale}
\newcommand{\wtBell}[3]{%
  \begin{scope}[shift={(#1,#2)},scale=#3]
    \draw[fishdetail] (0,0) circle (0.06);
    \fill[telosink] (-0.06,-0.06) .. controls (-0.22,-0.20) and (-0.26,-0.48) ..
      (-0.24,-0.66) -- (0.24,-0.66) .. controls (0.26,-0.48) and (0.22,-0.20) ..
      (0.06,-0.06) -- cycle;
    \draw[fishdetail] (-0.24,-0.66) -- (0.24,-0.66);
  \end{scope}}

% Cylinder weight for a drop shot, with its pinch-clip top.
\newcommand{\wtCylinder}[3]{%
  \begin{scope}[shift={(#1,#2)},scale=#3]
    \draw[fishheavy] (0,0) -- (0,-0.12);
    \draw[fishheavy] (-0.07,-0.12) -- (0.07,-0.12);
    \fill[telosink] (-0.10,-0.16) rectangle (0.10,-0.72);
    \fill[telosink] (0,-0.72) ellipse (0.10 and 0.06);
  \end{scope}}

% Split shot: a sphere with the pinch slot showing.
\newcommand{\wtSplitShot}[3]{%
  \begin{scope}[shift={(#1,#2)},scale=#3]
    \fill[telosink] (0,0) circle (0.15);
    \draw[telospaper,line width=0.5pt] (-0.17,0) -- (0.17,0);
  \end{scope}}

% ---------------------------------------------------------------------------
% Hardware
% ---------------------------------------------------------------------------

% Barrel swivel, vertical. \swivel{x}{y}{scale}
\newcommand{\swivel}[3]{%
  \begin{scope}[shift={(#1,#2)},scale=#3]
    \draw[fishheavy] (0,0.26) circle (0.09);
    \draw[fishheavy] (0,-0.26) circle (0.09);
    \fill[telosink] (-0.09,-0.17) rectangle (0.09,0.17);
    \draw[telospaper,line width=0.4pt] (-0.09,0) -- (0.09,0);
  \end{scope}}

% Three-way swivel: three eyes on a central ring.
\newcommand{\swivelThree}[3]{%
  \begin{scope}[shift={(#1,#2)},scale=#3]
    \fill[telosink] (0,0) circle (0.10);
    \foreach \a in {90,-30,210} {
      \draw[fishheavy] (\a:0.10) -- (\a:0.22);
      \draw[fishheavy] (\a:0.30) circle (0.085);}
  \end{scope}}

% Bead. \bead{x}{y}{scale}
\newcommand{\bead}[3]{%
  \begin{scope}[shift={(#1,#2)},scale=#3]
    \draw[fishheavy] (0,0) circle (0.15);
    \draw[fishdetail] (0,0) ellipse (0.05 and 0.15);
  \end{scope}}

% Bobber stop knot on the line.
\newcommand{\stopknot}[3]{%
  \begin{scope}[shift={(#1,#2)},scale=#3]
    \foreach \i in {-1,0,1} {
      \draw[fishheavy] (\i*0.075,-0.13) .. controls (\i*0.075+0.06,0) .. (\i*0.075,0.13);}
    \draw[fishheavy] (-0.19,0.16) -- (-0.05,0.10);
    \draw[fishheavy] (0.19,0.16) -- (0.05,0.10);
  \end{scope}}

% Slip float, sitting in the water: lower half filled.
\newcommand{\slipfloat}[3]{%
  \begin{scope}[shift={(#1,#2)},scale=#3]
    \draw[fishheavy] (0,0) ellipse (0.26 and 0.50);
    \begin{scope}
      \clip (-0.30,-0.56) rectangle (0.30,0);
      \fill[telosink] (0,0) ellipse (0.26 and 0.50);
    \end{scope}
    \draw[fishheavy] (0,0.50) -- (0,0.88);
    \draw[fishheavy] (0,-0.50) -- (0,-0.76);
    \draw[fishdetail] (-0.05,0.86) -- (0.05,0.86);
  \end{scope}}

% ---------------------------------------------------------------------------
% Baits
% ---------------------------------------------------------------------------

% Soft-plastic stick worm, drawn hanging and slightly curved.
\newcommand{\baitWorm}[3]{%
  \begin{scope}[shift={(#1,#2)},scale=#3]
    \def\wormbody{
      (0.02,0.06) .. controls (0.26,-0.10) and (0.20,-0.60) .. (0.06,-1.02)
      .. controls (-0.02,-1.28) and (-0.14,-1.34) .. (-0.20,-1.20)
      .. controls (-0.26,-1.06) and (-0.14,-0.98) .. (-0.14,-0.86)
      .. controls (-0.14,-0.50) and (-0.24,-0.14) .. (-0.14,0.02)
      .. controls (-0.10,0.10) and (-0.04,0.10) .. (0.02,0.06)}
    \fill[telospaper] \wormbody -- cycle;
    \draw[fishheavy] \wormbody -- cycle;
    \foreach \y/\a/\b in {-0.18/-0.19/0.16, -0.42/-0.20/0.19, -0.66/-0.17/0.18,
                          -0.88/-0.13/0.13} {
      \draw[fishdetail] (\a,\y) -- (\b,\y);}
  \end{scope}}

% Live minnow, side on, nose at the anchor and tail to the right.
\newcommand{\baitMinnow}[3]{%
  \begin{scope}[shift={(#1,#2)},scale=#3]
    \fill[telospaper]
      (0,0) .. controls (0.30,0.26) and (0.78,0.24) .. (1.02,0.08)
      .. controls (0.78,-0.24) and (0.30,-0.26) .. (0,0) -- cycle;
    \draw[fishheavy]
      (0,0) .. controls (0.30,0.26) and (0.78,0.24) .. (1.02,0.08)
      .. controls (0.78,-0.24) and (0.30,-0.26) .. (0,0) -- cycle;
    \draw[fishheavy] (1.02,0.08) -- (1.30,0.26) -- (1.26,-0.06) -- (1.02,0.08) -- cycle;
    \draw[fishdetail] (0.22,0.10) .. controls (0.30,0.02) and (0.30,-0.02) .. (0.22,-0.10);
    \fill[telosink] (0.13,0.05) circle (0.035);
    \draw[fishdetail] (0.44,0.20) -- (0.62,0.14);
    \draw[fishdetail] (0.40,-0.13) .. controls (0.50,-0.20) and (0.58,-0.20) .. (0.64,-0.14);
  \end{scope}}

% Nightcrawler threaded on a hook, with a free tail.
\newcommand{\baitCrawler}[3]{%
  \begin{scope}[shift={(#1,#2)},scale=#3]
    \draw[fishheavy] (0,0) .. controls (0.16,-0.34) and (-0.14,-0.56) ..
      (0.02,-0.90) .. controls (0.14,-1.14) and (-0.10,-1.24) .. (-0.16,-1.44);
    \foreach \i in {1,...,7} {
      \draw[fishdetail] ($(-0.06,{-0.14*\i})$) -- ($(0.06,{-0.14*\i-0.02})$);}
  \end{scope}}

% Leech, drawn with its characteristic ribbon swim.
\newcommand{\baitLeech}[3]{%
  \begin{scope}[shift={(#1,#2)},scale=#3]
    \def\leechbody{
      (0,0.02) .. controls (0.30,0.26) and (0.64,-0.22) .. (0.98,0.06)
      .. controls (0.94,0.18) and (0.86,0.20) .. (0.80,0.14)
      .. controls (0.56,-0.06) and (0.36,0.30) .. (0.10,0.16)
      .. controls (0.02,0.12) and (-0.02,0.08) .. (0,0.02)}
    \fill[telospaper] \leechbody -- cycle;
    \draw[fishheavy] \leechbody -- cycle;
  \end{scope}}

% ---------------------------------------------------------------------------
% Frame furniture for a rig plate
% ---------------------------------------------------------------------------

% Water surface: a ruled line with a wave band beneath it.
\newcommand{\rigsurface}[3]{%
  \draw[fishdetail] (#1,#3) -- (#2,#3);
  \draw[telosrule,line width=0.4pt,decorate,
        decoration={snake,amplitude=0.5mm,segment length=3.4mm}] (#1,#3-0.13) -- (#2,#3-0.13);
  \node[riglabel,anchor=west,text=telosmute] at (#1+0.04,#3+0.20) {surface};}

% Lake bottom: a ruled line with hatching below it.
\newcommand{\rigbottom}[3]{%
  \draw[fishheavy] (#1,#3) -- (#2,#3);
  \foreach \x in {#1,...,#2} {
    \draw[fishdetail] (\x+0.5,#3) -- (\x+0.26,#3-0.24);}
  \node[riglabel,anchor=west,text=telosmute] at (#1+0.04,#3-0.34) {bottom};}

% A labelled callout with a hairline leader and a dot at the component end.
%   \callout{from x}{from y}{to x}{to y}{anchor}{text}
\newcommand{\callout}[6]{%
  \draw[fishdetail,telosmute] (#1,#2) -- (#3,#4);
  \fill[telosmute] (#1,#2) circle (0.035);
  \node[riglabel,anchor=#5,inner sep=1.6pt] at (#3,#4) {#6};}

% A dimension arrow with a label, for leader lengths and set depths.
\newcommand{\rigspan}[5]{%
  \draw[|<->|,fishdetail,telosmute] (#1,#2) -- (#3,#4);
  \node[rigtag,anchor=#5,inner sep=2pt,fill=telospaper]
    at ({(#1+#3)/2},{(#2+#4)/2}) {}; }
\newcommand{\rigspanlabel}[6]{%
  \draw[|<->|,fishdetail,telosmute] (#1,#2) -- (#3,#4);
  \node[rigtag,inner sep=2pt,fill=telospaper,anchor=#5]
    at ({(#1+#3)/2},{(#2+#4)/2}) {#6};}


% Cooking sheets: filleting overlays and the consumption advisory.
% Shared machinery for the cooking sheets.
%
% The filleting diagrams are drawn *on top of* the same public-domain Denton
% plate the species sheet uses. That is deliberate: a reader who has just
% identified a fish from one sheet sees the identical drawing on the other, with
% the cuts marked on it, so there is never any doubt which fish or which end.
%
% Cut lines are drawn twice -- a white halo first, then the black dashed line on
% top -- because a plain black line disappears into the dark dorsal shading of a
% grayscale plate.

\tikzset{
  cuthalo/.style={draw=telospaper,line width=2.6pt,line cap=round},
  cut/.style={draw=telosink,line width=1.1pt,dash pattern=on 3.2pt off 2.2pt,
              line cap=round},
  cutsolid/.style={draw=telosink,line width=1.1pt,line cap=round},
  cutguide/.style={draw=telosmute,line width=0.5pt,dash pattern=on 1.2pt off 1.2pt},
  cutlabel/.style={font=\fontsize{6.4}{7.4}\selectfont},
}

% The plate, set up so overlay coordinates run 0..1 across and up the image.
%   \begin{filletplate}{width}{plate-slug} ...overlay... \end{filletplate}
\newenvironment{filletplate}[2]
  {\begin{tikzpicture}[x=#1,y=#1]
     \node[anchor=south west,inner sep=0] (telosfishplate) at (0,0)
       {\includegraphics[width=#1]{\telosplateroot/#2.png}};
   \begin{scope}[x={(telosfishplate.south east)},y={(telosfishplate.north west)}]}
  {\end{scope}\end{tikzpicture}}

% A cut line: halo then dashes. Takes a coordinate path fragment.
\newcommand{\cutpath}[1]{%
  \draw[cuthalo] #1;
  \draw[cut] #1;}

% A solid line, for a cut that is a single decisive stroke rather than a
% traced-around outline.
\newcommand{\cutsolidpath}[1]{%
  \draw[cuthalo] #1;
  \draw[cutsolid] #1;}

% Numbered step marker, matching the numbered spots on the lake maps so the
% guide has one visual language for "this is step/place N".
\newcommand{\cutstep}[3]{%
  \node[circle,fill=telosink,text=telospaper,inner sep=0pt,minimum size=3.4mm,
        font=\fontsize{6.4}{6.4}\selectfont\bfseries] at (#1,#2) {#3};}

% A small caption pinned inside the plate area.
\newcommand{\cutnote}[4]{%
  \node[cutlabel,anchor=#3,fill=telospaper,fill opacity=0.82,text opacity=1,
        inner sep=1.4pt] at (#1,#2) {#4};}

% ---------------------------------------------------------------------------
% The eating-quality strip. Same shape on every cooking sheet so the sheets can
% be compared by eye: is this fish worth carrying home?
% ---------------------------------------------------------------------------
\newenvironment{eatingfacts}
  {\begingroup\small
   \setlength{\parskip}{0pt}%
   \renewcommand{\arraystretch}{1.22}%
   \begin{tabular}{@{}>{\bfseries\raggedright\arraybackslash}p{0.19\linewidth}@{\hspace{0.6em}}>{\raggedright\arraybackslash}p{\dimexpr0.81\linewidth-0.6em\relax}@{}}}
  {\end{tabular}\par\endgroup}

% ---------------------------------------------------------------------------
% The Wisconsin advisory. Written once, verbatim from the DNR/DHS "Choose
% Wisely" 2024--2026 statewide guidance, so no sheet can paraphrase it loosely.
% The argument is the row that applies to this species.
% ---------------------------------------------------------------------------
\newcommand{\advisorystrip}[3]{%
  \par\vspace{0.35em}
  \noindent
  \begin{minipage}{\linewidth}
    \rule{\linewidth}{0.9pt}\par
    \vspace{0.3em}
    {\footnotesize\textbf{Wisconsin safe-eating guidance.} For #1:
     \textbf{#2} for women under 50 and children under 15;
     \textbf{#3} for women over 50 and men.
     This is the statewide advice for inland waters, from the DNR and DHS
     \emph{Choose Wisely} guide. Over 150 Wisconsin waters carry stricter
     site-specific advice --- check the current booklet or the DNR's
     Eat Your Catch query tool for these lakes before you make a habit of it.\par}
    \vspace{0.25em}
    \rule{\linewidth}{0.3pt}
  \end{minipage}\par}

% The three statewide categories, so a sheet names one rather than retyping it.
\newcommand{\advisorypanfish}{\advisorystrip{bluegill, crappie, yellow perch,
  sunfish, rock bass and bullhead}{one serving per week}{unrestricted}}
\newcommand{\advisorygamefish}[1]{\advisorystrip{#1}{one serving per month}%
  {one serving per week}}
\newcommand{\advisorymusky}{\advisorystrip{muskellunge}{do not eat}%
  {one serving per month}}

